% ---------------------------------------------------------------------------
% Response to R2 Major 2 -- JEDC revision
% Compile with pdfLaTeX. No figures required; standard packages only.
% To paste into an existing response document, take everything between the
% BODY BEGINS and BODY ENDS markers.
% ---------------------------------------------------------------------------
\documentclass[11pt]{article}

\usepackage[utf8]{inputenc}
\usepackage[T1]{fontenc}
\usepackage{lmodern}
\usepackage[margin=1in]{geometry}
\usepackage{booktabs}
\usepackage{amsmath}
\usepackage{microtype}
\usepackage{framed}
\usepackage{xcolor}
\usepackage[hidelinks]{hyperref}
\colorlet{shadecolor}{gray!12}

\setlength{\parskip}{0.6em}
\setlength{\parindent}{0pt}

\title{\textbf{R2, Major 2}\\[0.3em]
       \large\normalfont\itshape Do the comparison economies face the same oil price?}
\author{}
\date{}

\begin{document}
\maketitle
\vspace{-3em}

% =========================== BODY BEGINS ===================================

The referee is right, and the objection was serious. In the submitted version
the constant-price economy held the oil price at the median of its log-normal
distribution rather than at its mean. The median is \textbf{1.1275}; the average
price the full economy actually pays is \textbf{1.2429}. So the economy we were
calling ``the same, without oil price risk'' was in fact buying oil
\textbf{10.2\% cheaper}, every quarter, for ever. Any welfare gain or spread
difference we attributed to removing risk included a permanent discount on oil.

The same problem, in smaller form, affected the constant-volatility economy.
Fixing volatility at the middle node of its grid is not the same as fixing it at
its average level: because the innovation variance enters exponentially, the
average of the variance exceeds the variance at the average, and an economy with
a more variable log price faces a higher average price level. This is the
convexity point the referee makes, and it ran in the same direction as the first
problem.

\section*{What we changed}

All four economies are now solved on one shared 277-node price grid, with the
same preferences, technology, default costs and shocks. Two pins moved:

\begin{itemize}
\item the constant-price economy is now held at $\mathrm{E}[p]$, the mean of the
      price level, rather than at the median of its logarithm;
\item the constant-volatility economy is now held at the average innovation
      variance, not at the middle node of the volatility grid.
\end{itemize}

The result is that the four economies face effectively the same oil price:

\begin{table}[h]
\centering
\small
\begin{tabular}{lcc}
\toprule
economy & average oil price in simulation & difference from the full model \\
\midrule
full                 & 1.2429 & --- \\
mean-preserving      & 1.2418 & $-0.10\%$ \\
constant volatility  & 1.2421 & $-0.07\%$ \\
constant price       & 1.2423 & $-0.05\%$ \\
\bottomrule
\end{tabular}
\end{table}

Within one tenth of one percent, against the 10\% gap in the submitted version.
Whatever now separates these economies is risk, not the price of oil.

\section*{What it does to the results}

With the average price held fixed, removing \emph{all} oil price risk is worth
between \textbf{0.34\% and 0.46\% of consumption}, depending on the level of
output, and removing only the \emph{time variation} in volatility between
\textbf{0.19\% and 0.22\%}. Both curves are positive everywhere and properly
nested, which is what removing strictly more risk has to deliver, and Figure~7
shows that.

These are smaller than the \textbf{1.1\%} in the submitted paper, but we should
not claim the whole of that difference for this correction. The substitution
elasticity has also been recalibrated since submission, in response to Major 3
of the other report, and the welfare cost of oil price risk is sensitive to it.
We have not re-run the matched comparison at the original elasticity, so we
cannot divide the change between the two. What we can say is the referee's point
itself: under the old construction the comparison economy bought oil 10.2\%
cheaper for ever, and that discount was being counted as a gain from removing
risk.

The spread comparison survives better. The difference in mean spreads between
the full economy and the constant-price economy is \textbf{200.7 basis points},
and because both now face the same average oil price this is a risk effect
rather than a price-level effect. It is about half the mean spread, so oil price
risk still matters a great deal to these borrowers.

The referee also asks for a comparison economy that varies only in dispersion.
We have built it, and it is the subject of our reply to Major 1.

\section*{Where this leaves the paper}

Section 6.2 and the welfare discussion have been rewritten around the matched
economies. The headline welfare number falls from 1.1\% to roughly 0.4\%, and we
say plainly in the text why it fell.

\bigskip
\hrule
\bigskip

\subsection*{Suggested manuscript wording}

\begin{shaded}
\noindent
The comparison economies are constructed to face the same average oil price as
the full model. The constant-price economy is pinned at $\mathrm{E}[p]$ rather
than at the median of the log price distribution, and the constant-volatility
economy at the average innovation variance rather than at the central volatility
node; simulated average prices then agree across the four economies to within
0.1\%. Under this construction, removing all oil price risk is worth between
0.34\% and 0.46\% of consumption and removing only the time variation in
volatility between 0.19\% and 0.22\%, against a mean-spread difference of 200.7
basis points for oil price risk as a whole. An earlier version of this
comparison held the constant-price economy at the median of the log-normal
distribution, so that it faced oil roughly 10\% cheaper on average; the welfare
gains reported there were correspondingly larger.
\end{shaded}

\subsection*{Note}

The four economies are \texttt{model/economies\_common\_check/\{benchmark,
mean\_preserving, intermediate, basic\}}, one standalone Fortran source each,
and the pins are the \texttt{p\_grid} assignments in the constant-price and
constant-volatility sources. Simulated average prices are computed over the
whole simulation, not over pre-default windows: the table above is
\texttt{table6\_mean\_price.csv}, written by
\texttt{model/analysis/JEDC\_decomposition.m} to
\texttt{model/economies\_common\_check/decomposition/} along with the other
tables and the figure. Welfare curves are from
\texttt{ALL\_FIGURES\_common\_grid.m}, Figure~7.

% ============================ BODY ENDS ====================================

\end{document}
